\section*{Support-twelve involution and square-root cover closure}
\addcontentsline{toc}{section}{Support-twelve involution and square-root cover closure}

\paragraph{Pass 4721: disjoint support turns the $1620$ minima into $540$ triangles.}
Let $\mathcal T$ be the exact support-$12$ minimum shell from Pass~4703, so
$|\mathcal T|=1620$ and every element is the corner-star thickening $T(A)$ of a
unique W33 apartment.  Join two thickenings when their twelve-line supports are
disjoint.  Exact enumeration gives degree two at every vertex and
\[
 \boxed{\Delta_{12}=540K_3.}
\]
The three supports in a component cover $36$ of the $40$ W33 lines.  Their
four-line complement is always pairwise disjoint.  The $540$ components yield
exactly $270$ distinct complements, each with two preimages.

\paragraph{Pass 4722: the old unnamed $270$ is the four-fixed-line involution orbit.}
The inner group $G=\PSp(4,3)$ has $315$ involutions in the explicit line action,
with fixed-line census
\[
 \boxed{4^{270}\oplus16^{45}.}
\]
The $270$ four-fixed-line sets are distinct and equal exactly to the $270$
Pass~4721 residues.  A residue stabilizer has order $96$, hence orbit size
$25920/96=270$, and each residue determines a unique inner involution.
This corrects the negative conclusion in historical Pass~1830: that GAP script
defined its permutation group on the $40$ W33 \emph{points} and then applied it
to a set of \emph{line indices}.  The reported orbit $2880$ was therefore not
the induced action on the four-line object.  Under the actual line action the
four-line invariant is the desired $270$-orbit.

\paragraph{Pass 4723: the two triangle sheets are outer order-four square-root sheets.}
Fix a residue $R$ and its unique inner involution $g$.  In
$\PGSp(4,3)$ there are exactly eight outer solutions of $h^2=g$; all have order
four, and exactly two fix four W33 lines.  Those two are $h$ and $h^{-1}$.
There are also exactly two support-disjointness triangles over $R$, and either
four-fixing root swaps the two.  More strongly, for either triangle $C$ over
$R$ and either four-fixing root,
\[
 \boxed{\operatorname{Stab}_{G}(C)=C_G(h),\qquad |C_G(h)|=48.}
\]
Thus
\[
 \boxed{
 \{540\text{ support-disjointness triangles}\}
 \cong_G
 \{540\text{ four-fixing outer order-four elements}\}.
 }
\]
Both are two-fold covers of the same $270$ residue/involution orbit.  The two
$G$-equivariant identifications differ by $h\leftrightarrow h^{-1}$.
The full $\PGSp(4,3)$ triangle stabilizer and root centralizer both have order
$96$ but are distinct extensions with common inner intersection of order $48$;
accordingly no untwisted full-$\PGSp$ identification is claimed.

\paragraph{Pass 4724: the $270$ residues factor the complement line relation.}
Let $B\in\{0,1\}^{40\times270}$ be line-versus-residue incidence and let $A_*$
be the W33 line-intersection adjacency matrix.  Every residue has four lines,
every line lies in $27$ residues, every skew line-pair lies in exactly three,
and an intersecting pair lies in none.  Hence
\[
 \boxed{BB^{\mathsf T}=27I+3(J-I-A_*).}
\]
Since the complement of $\operatorname{SRG}(40,12,2,4)$ has spectrum
$27^1\oplus(-3)^{24}\oplus3^{15}$,
\[
 \boxed{\operatorname{spec}(BB^{\mathsf T})
 =108^1\oplus18^{24}\oplus36^{15}.}
\]
Thus $B$ has real rank $40$.  Over $\mathbb F_2$ its residue-mask span has rank
$30$; no module identification is inferred from that dimension alone.  A
representative $270$-class involution fixes exactly $24$ support-$12$
thickenings and eight of the $540$ triangles.

\begin{center}
\fbox{\parbox{0.94\linewidth}{\textbf{Evidence boundary.}
These are exact finite incidence and permutation-group statements verified by
\texttt{analysis/w33\_pass4721\_4724\_support12\_involution\_square\_root\_cover.py}.
The result repairs an action-domain error in Pass~1830 and then adds a genuine
minimum-shell/square-root-cover theorem.  Equal cardinalities are never used as
a substitute for an intertwiner: fixed sets, square maps, stabilizers,
centralizers, orbits, and the Gram identity are checked explicitly.  No physical
interpretation is inferred.}}
\end{center}
